\section{Discussion}

This section discusses the main findings of the proposed pipeline and their implications for preoperative assessment in root canal treatment, and then outlines the limitations of the current work, including the sample and data source, the choice of reference standard, and the evaluation design.

\subsection{Main Findings}
Without any physical calibration object, the proposed pipeline reduced the filling-length estimation error to about 1~mm. Compared with unprocessed images, top-hat preprocessing reduced the MAE from 1.433 to 0.993~mm and increased $R$ from 0.232 to 0.727. These findings indicate that length-estimation performance depends not only on the segmentation network but also on the quality of the input image; in particular, suppressing large-scale exposure non-uniformity and enhancing local contrast at the root--bone boundary reduces uncertainty at the apical boundary and stabilises the pixel length measured along the principal axis.

Long-side normalisation and the safety offset address scale and clinical safety, respectively. Using only the image size, long-side normalisation reduces the between-image variation of px/mm and markedly increases the correlation between pixel length and true length. The fixed 0.5-mm offset keeps the overall bias at $-0.493$~mm so that estimates tend to be short, consistent with the clinical principle of avoiding overextension and reducing the risk of instruments or filling material passing beyond the apical foramen. Regarding the regression model, GPR and the ANN performed similarly, but GPR requires no random initialisation, is deterministic, and outperformed linear regression with the same inputs, making it suitable for the small-sample setting of this study.

\subsection{Limitations}
Although the proposed pipeline achieved favourable accuracy, the current evaluation remains preliminary. First, only 67 anterior teeth with reference filling lengths were tested, all from a single centre, and posterior and multi-rooted teeth were not included; the choices of preprocessing and input features were also based on cross-validation results on the same 67 teeth, which may make the reported performance slightly optimistic. Second, the post-treatment filling length was used as a surrogate for the working length, and the two are not identical in clinical practice; moreover, a periapical radiograph is a two-dimensional projection and cannot fully represent the true length of curved canals.

In addition, the target tooth in each test image was manually selected after detector cropping, so the pipeline is not yet fully automatic at this step; tooth contours were annotated on CLAHE images and reused for all preprocessing groups, so the annotations themselves may be biased towards that preprocessing, and the top-hat segmentation model was trained with a single random seed, so the effect of training stochasticity has not been quantified. Future studies should include more tooth types and samples, multicentre data, independent annotation by multiple clinicians with inter-observer agreement analysis, and prospective clinical comparison with current electronic apex locators and intraoperative radiographs.
